


International migration of Nepalese people has been a common phenomenon even before nineteenth century when they used to travel to Lahore to join the army of the Sikh ruler Ranjit Singh which was further institutionalised with the recruitment of Nepalese to the British Gurkhas in 1815-16. \cite{Seddon2002} Foreign migration has been a by-product of stagnant rural economy 

\cite{Voors2012} advocate that there are two pathways for an economy after civil war; if it destroys the social capital or increases impatience then it will reduce investment and growth, however if it triggers institutional improvements or alters preferences towards saving then income might even be higher than before. 